\documentclass[12pt]{article}
\usepackage[utf8]{inputenc}
\usepackage[T1]{fontenc}
\usepackage{hyperref}
\title{Koroni Language Resources for K-12 Curriculum Development}
\author{Compiled by Claude Sonnet 5.5 (claude-sonnet-5-5) for Andr\'as Kornai}
\date{}
\begin{document}
\maketitle

\section{Introduction}
Koroni (ISO 639-3 xkq; Glottocode koro1311) is an Austronesian language of Sulawesi, Indonesia, in the East Coast Bungku-Tolaki group of Celebic. Verbix places it in Unsongi village, Bungku Tengah subdistrict, on the east coast of Tolo Bay. It is also called Oengsongi. Speaker estimates are rough: about 500 (Verbix) and about 1,000 for the people group (Joshua Project). Wikipedia describes Kulisusu, Koroni and Taloki as a dialect chain, and Verbix gives lexical similarity of 75\% with Taloki and Kulisusu and 65\% with Moronene.

\section{Primary Educational Resources}
No primers, textbooks or school programmes for Koroni were found. Joshua Project notes that written or published materials are limited or uncertain.

\section{Linguistic and Reference Materials}
\begin{itemize}
\item Glottolog, languoid koro1311. \url{https://glottolog.org/resource/languoid/id/koro1311.bigmap.html}
\item Verbix, Koroni language page. \url{https://docs.verbix.com/Languages/Koroni}
\item Wikipedia, ``Kulisusu language'': describes the Kulisusu--Koroni--Taloki dialect chain, so related-language materials may help. \url{https://en.wikipedia.org/wiki/Kulisusu_language}
\item Endangered Languages Project sample page (listed in search results; it may concern a different language, so check before use). \url{https://endangeredlanguages.com/lang/1902/samples/11630}
\end{itemize}

\section{Religious and Community Resources}
\begin{itemize}
\item Joshua Project, ``Koroni in Indonesia'': people-group profile. \url{https://joshuaproject.net/people_groups/12801}
\item Joshua Project, language page for xkq. \url{https://joshuaproject.net/languages/xkq}
\end{itemize}

\section{Digital and Audio Resources}
No audio, video or corpus resources specific to Koroni were confirmed.

\section{Curriculum Development Context}
Documentation is minimal. A curriculum would need local documentation work, an orthography agreed with the community, and coordination with Indonesian schooling, where local-language content is possible. Related Kulisusu and Taloki materials, if they exist, might serve as a starting point, but they have not been checked here.

\end{document}